\subsection{四単純位数4の16次元atlasと第一rank型の完成}
\label{subsec:cubic-four-simple-order4-rankC}
以下では中山スペクトルを $(1,-1,i,-i)$、固有基底を $(x,y,b,c)$ とする。
これは系\ref{cor:cubic-detminus1-spectral-boundary}の有限残余を含むが、
全四単純スペクトルをこの一点へ制限する主張ではない。
重みは $\mathbb Z/4$ で $(0,2,1,3)$ とする。

\begin{proposition}[元16tensorと三つの必要rank型]
\label{prop:cubic-four-simple-order4-atlas}
このスペクトルの実際の中山巡回式を満たす四次最高tensor空間は16次元である。
二次の重み $0,2,1,3$ の順に関係数を書けば、H3に必要なrank型は
\[
 (2,2,1,1),\qquad (2,0,2,2),\qquad (0,2,2,2)
\]
のみである。どの型でも $B_3$ の各重みの次元は5となる。
これらは必要条件であり、AS4又はH3十分性ではない。
\end{proposition}
\begin{proof}
全四次語を右巡回軌道に分け、回転時に末字の中山固有値の逆数の負を掛ける。
一周した係数が1となる軌道の代表は、順に
\[
\begin{gathered}
xxyy,xxbc,xxcb,xybb,xycc,xbxc,xbyb,xbby,\\
xcyc,xccy,yyyy,yybc,yycb,ybyc,bbcc,bcbc.
\end{gathered}
\]
その係数を $a_0,\ldots,a_{15}$ とする。
二次の基底を
\[
\begin{array}{c|l}
0&xx,yy,bc,cb\\2&xy,yx,bb,cc\\1&xb,yc,bx,cy\\3&xc,yb,by,cx
\end{array}
\]
と取る。middleの自己blockとpaired blockは正確に
\[
 H_0=\begin{pmatrix}
0&a_0&a_1&a_2\\a_0&a_{10}&a_{11}&a_{12}\\
a_1&a_{11}&a_{15}&-ia_{14}\\a_2&a_{12}&-ia_{14}&-ia_{15}
\end{pmatrix},\quad
 H_2=\begin{pmatrix}
0&-a_0&a_3&a_4\\a_0&0&a_7&a_9\\
-a_3&-a_7&0&a_{14}\\-a_4&-a_9&-a_{14}&0
\end{pmatrix},
\]
\[
 M=\begin{pmatrix}
a_5&a_6&a_7&-a_1\\-ia_8&-ia_{13}&a_{12}&-a_4\\
ia_2&ia_3&ia_6&-a_5\\-ia_9&-ia_{11}&-ia_{13}&-a_8
\end{pmatrix}.
\]
従って $r_2$ は偶数、$r_1=r_3$、$r_0+r_2+2r_1=6$。
同変Koszul解消とtopのhdetが1であることから、
$\mu$ と $\mu^2$ のtrace級数の分母はそれぞれ
\[
 1+(r_0-r_2)t^2+t^4,\qquad 1+(6-4r_1)t^2+t^4.
\]
有限位数なのでtrace係数はHilbert値で絶対値を抑えられる。
分母は単位円内部に零を持たず、相反性から全根は単位円上にある。
$1+dt^2+t^4$ がこの条件を満たすには $|d|\le2$ が必要である。
従って $|r_0-r_2|\le2,\ |6-4r_1|\le2$。
関係数6と偶数性を合わせると表示した三型だけとなる。
最後に $V=(1,1,1,1)$ を使い、character恒等式
$[B_3]=V^3-2VI+V^*$ にいずれの型を代入しても $(5,5,5,5)$ を得る。
\end{proof}

\begin{lemma}[自己拡大上で自明なtwisted Serre作用の障害]
\label{lem:cubic-twisted-character-Serre}
Hom有限な三角圏で $L$ は単純対象、$S$ はSerre関手とする。
$S(L)\simeq L[2]$ を固定し、これによる $S[-2]$ の
$\Ext^1(L,L)$ 上の作用が恒等であるとする。
この自己拡大空間が一次元なら矛盾する。
\end{lemma}
\begin{proof}
Serre traceとshiftの符号から、任意の自己一次拡大 $a,b$ について
\[
 \operatorname{tr}(b[1]a)=-\operatorname{tr}((S[-2]a)[1]b)
\]
となる。作用が恒等ならこの対は交代である。
Serre双対性から非退化でなければならないが、標数0の一次元空間には
非退化交代形式は存在しない。
ここでは $S(L)\simeq L[2]$ だけから作用の自明性を推論していない。
\end{proof}

\begin{theorem}[重み0関係が零の全rank型の排除]
\label{thm:cubic-four-simple-order4-rankC-empty}
上の $(0,2,2,2)$ 型は全てH3を与えない。
\end{theorem}
\begin{proof}
$H_0=0$ なので $a_0,a_1,a_2,a_{10},a_{11},a_{12},a_{14},a_{15}=0$。
$H_2$ のrank2は
\[
 \operatorname{rank}\begin{pmatrix}a_3&a_4\\a_7&a_9\end{pmatrix}=1
\]
と同値である。その像には純 $xy,yx$ 関係と純 $b^2,c^2$ 関係が一本ずつある。
整域性から $a_3,a_4,a_7,a_9$ は全て非零。
対角generator/top変更で $a_3=a_4=1,a_7=a_9=q\ne0$ とできる。
残る係数を $A=a_5,h=a_6,k=a_8,E=a_{13}$ とする。
$M$ の $3\times3$ minorsのGroebner基底には
\[
\begin{gathered}
 A(q^4-1),\ E(q^4-1),\ h^2-k^2,\ k^3-qk,\\
 k^2(q^2+1)-q(q^2+1),\ AE+qk^2-q^2,\\
 A(h-qk),\ E(h+qk),\
 A\bigl(k^2-(q^3+q)/2\bigr),\
 E\bigl(k^2-(q^3+q)/2\bigr)
\end{gathered}
\]
が含まれる。全13本と逆向きのideal同定は公開再現器に保持する。
このpaired blockにはminor $-q$ があるので、rankは全値で正確に2である。

まず $A=E=0$ の枝では $k^2=q,h=\epsilon k$、$\epsilon=\pm1$ となる。
$k=t\ne0$ と書くと、全値の元関係画像は
\begin{gather*}
 b^2+c^2=0,\quad xy+t^2yx=0,\\
 xb=-i\epsilon t^{-1}bx,\quad xc=it^{-1}cx,\quad
 yb=-\epsilon tby,\quad yc=-tcy.
\end{gather*}
これは二元AS2平面 $T(b,c)/(b^2+c^2)$ 上の逐次skew Ore拡大である。
対角自己同型は関係を保ち互いに可換であるから、全 $t\ne0$ でこの表示とPBWが成立する。
元middle画像との同定はunit六次minorとLaurent-$t$ 座標によっても検算する。
固定三次の基底は
\[
 x^3,\quad xy^2,\quad xbc,\quad xcb,\quad yb^2.
\]
$x,y$ の正常化自己同型のこの基底上の固有値は、$q=t^2$ としてそれぞれ
\[
 (1,q^2,\epsilon/q,\epsilon/q,1),\qquad
 (-q^{-3},-q^{-1},-\epsilon,-\epsilon,q).
\]
$Gx=xG,Gy=-yG$ が必要なので、$q^3=1$ なら全候補は左因子 $x$ を持ち、
$q=-1$ なら $kyb^2$ に含まれ、それ以外では零となる。
特に $xbc,xcb$ が残るのは $q=1,\epsilon=1$ だけであり、やはり因子 $x$ を持つ。
いずれも三次商は整域でない。

$A,E$ のいずれかが非零なら $q^4=1$。
上の全式を使うと、残余は次の十個の対角軌道で覆われる。
$x\mapsto ux,y\mapsto u^{-1}y$ は $A,E$ を $u^2A,u^{-2}E$ 倍するので、
一方だけ非零ならそれを1へ、両方非零なら $A=1$ へできる。
\[
\begin{array}{c|c|c|c|c|c}
 q&A&E&k&h&\text{正常核又は障害}\\\hline
 1&1&0&\pm1&k&\dim2,\ \text{全て左因子 }x\\
 1&0&1&\pm1&-k&0\\
 -1&1&0&i&-i&0\\
 -1&1&0&-i&i&\dim1\\
 -1&0&1&\pm i&k&0\\
 \pm i&1&-1&0&0&0
\end{array}
\]
各軌道で元の全次数3,4商を構成し、固定三次5列の正常核をfreshに計算した。
表の左因子も $xB_2$ への全核の包含として厳密検算した。

最後の一次元正常核の包絡関係は
\begin{gather*}
 b^2+c^2=0,\quad xy=yx,\quad xb=-bx,\quad xc=-cx,\\
 yb=-iby+ixc,\qquad yc=icy+xb,
\end{gather*}
であり、その核は
\[
 G=c^2y-\tfrac12bcx+\tfrac i2cbx
\]
で張られる。$x$ は正常で、その共役は $y$ を固定し $b,c$ を負にする。
$R=B/(G)$ がH3だと仮定する。
$x$ のOre開集合の次数零環は、$Y=yx^{-1},B_0=bx^{-1},C_0=cx^{-1}$ により
\begin{gather*}
 B_0^2+C_0^2=0,\\
 YB_0-iB_0Y=iC_0,\qquad YC_0+iC_0Y=B_0,\\
 2C_0^2Y-B_0C_0+iC_0B_0=0
\end{gather*}
と表示される。
正常 $x$ のskew Laurent表示を両方向に構成すれば、このdehomogenizationは全関係を与える。
原点characterの自己一次拡大は表示関係の線型化の核であり、
第一、第四式の一次項は零、第二、第三式はそれぞれ $C_0,B_0$ を殺す。
従って $\Ext^1$ は $Y$ 方向の一次元である。

H3の $S=(-1)[2]$ をこの開集合へ制限し、$x$ で次数shiftを同定すると、
$S[-2]$ は $Y\mapsto Y,B_0\mapsto-B_0,C_0\mapsto-C_0$ で与えられる。
原点characterを固定し、自己拡大の唯一の $Y$ 方向にも恒等に作用する。
この点支持対象のExtとSerre traceはOre開集合で計算できる。
補題\ref{lem:cubic-twisted-character-Serre}に反し、最後の軌道もH3でない。
これでこのrank型の全必要被覆が空となる。
\end{proof}

公開再現器は \texttt{verify\_cubic\_four\_simple\_order4\_rankC.py}。
元16tensor、全rank算術、Groebner基底、全値source画像、十軌道の全商・正常核と因子、
例外の唯一の三次とcharacter線型化をexactに検算する。
PBW、因子による非整域性、twisted Serre traceの通常証明を計算成功だけで代替してはいない。
残る $(2,0,2,2),(2,2,1,1)$ と全四単純・全CY被覆は未完成である。
